\documentclass[11pt]{article}
\usepackage{amsfonts}
\usepackage{amsmath}
\usepackage{amssymb}
\usepackage{color}
\usepackage{fullpage}
\usepackage{graphicx}
\usepackage[colorlinks=true,urlcolor=darkgray,bookmarks=false]{hyperref}
\usepackage{multirow}
\usepackage{setspace}

\newtheorem{theorem}{Theorem}
\newtheorem{assumption}[theorem]{Assumption}
\newtheorem{axiom}[theorem]{Axiom}
\newtheorem{claim}[theorem]{Claim}
\newtheorem{conjecture}[theorem]{Conjecture}
\newtheorem{corollary}[theorem]{Corollary}
\newtheorem{definition}[theorem]{Definition}
\newtheorem{example}[theorem]{Example}
\newtheorem{lemma}[theorem]{Lemma}
\newtheorem{proposition}[theorem]{Proposition}
\newtheorem{remark}[theorem]{Remark}
\newcounter{question}
\newtheorem{problem}{Problem}[question]
\newenvironment{proof}[1][Proof]{\noindent\textbf{#1.} }{\ \rule{0.5em}{0.5em}}
\setlength{\parindent}{0mm}
\def\changemargin#1#2{\list{}{\rightmargin#2\leftmargin#1}\item[]}
\let\endchangemargin=\endlist
\definecolor{scarlet}{cmyk}{0,1.00,0.65,0.15}
\definecolor{gray}{cmyk}{0.06,0,0,0.34}
\definecolor{darkgray}{cmyk}{0.07,0,0,0.90}

\begin{document}

\thispagestyle{empty}
\vspace*{-20mm}

%%%% HEADER TABLE %%%%%%

\begin{changemargin}{-0.5in}{0.5in}
\footnotesize{
\begin{tabular}{lr}
\multirow{1}{*}{\includegraphics[height=0.66in,width=0.66in]{graphics/osu_logo.pdf}} & \vspace{-5pt}\textbf{\small{Department of Economics}}\\
&\color{scarlet}{\rule{6.375in}{0.4pt}}\hspace*{1pt}\\
& 410 Arps Hall\\
& 1945 North High Street\\
& Columbus, OH 43210-1172\\
& \\
& \textbf{Paul J. Healy}\\
& Professor of Economics\\
& Phone: (614) 247-8876\\
& Fax: (614) 292-3906\\
& Email: \href{mailto:healy.52@osu.edu}{healy.52@osu.edu}\\
& Web: \href{http://healy.econ.ohio-state.edu}{healy.econ.ohio-state.edu}
\end{tabular}
}
\end{changemargin}

%%%%%%  LETTER BODY  %%%%%%%%%

%\today

\vspace{1cm}

\begin{center}
    \Large{\textbf{Course Proposal:\\ ECON 8877: Experimental Economics Methodology}}
\end{center}

\bigskip
\setlength{\parskip}{10pt}

I propose adding a new course to the Economics course catalogue, ECON 8817, entitled ``Experimental Economics Methodology''. This would be added to the requirements for the Experimental Economics 2nd-year field sequence.

\section{Course Description}

\textbf{ECON 8877 - Experimental Economics Methodology}

\begin{description}
\item[Units:] 3 units
\item[Campus:] Columbus
\item[Grading Basis:] Graded A--E
\item[Career:] Graduate
\item[Course Components:] Lecture, required
\end{description}
\paragraph{Description} This course covers the key aspects of designing, running, and analyzing an incentivized experiment with human subjects. Topics include a theoretical analysis of incentives in experiments, incentive compatible elicitation techniques, statistical methodology, model selection, programming techniques, and ethical considerations.\\
Prereq: 8714 (808) or equiv with permission of instructor.

\section{Discussion of Impacts of Adding ECON 8877}

Currently, the Experimental Economics field offers only two courses: ECON 8873 - Industrial Organization III: Economic Analysis of Auctions, and ECON 8875 - Experimental and Behavioral Econoimcs. Students who opt for the Experimental Economics field are not currently required to take a 3rd class, though they are encouraged to take ECON 8817 - Game Theory (part of the Theory field sequence), a math course, or another course matching their interests. In practice, most students who take the Experimental Economics field also take the Theory field, in which case they must take ECON 8817 - Game Theory as part of the Theory field requirement. But, regardless of their second field, most students who opt for the Experimental field only take five classes in total. Adding ECON 8877 will simply add a sixth course for these students.

Will the addition of ECON 8877 as part of the Experimental field take away enrollments from other courses in Economics? For students who sign up for both the Theory and Experimental fields there should be no impact, since they must already take ECON 8817 - Game Theory as part of their Theory field requirement. This simply adds a sixth course to their requirements, increasing the total number of courses they take. Where there could be an impact is in the cases where a student signs up for the Experimental field and some other field, would have taken a different 6th course inside Economics, but now must take ECON 8877 instead. From our limited experience this seems to be a rare case (often the 6th course is taken outside of Economics), but is possible.

As for my own teaching, I've been on a two-course load for several years due to my service as Director of Graduate Admissions---and had a teaching reduction in 2022--23 thanks to a one-semester sabbatical---so I need to add a new class starting in 2023--24. Creating ECON 8877 at this time enables me to add a new class without having to switch away from teaching any other courses.

Is this course viable, given that the long-run viability of the Experimental field is a bit uncertain? Although we hope to hire new faculty in this area, it's possible we also experience retirements that would force us to reduce the field back to two courses, or possibly cancel the field altogether. However, in the case of retirements, it is likely we would cancel either ECON 8873 or ECON 8875 rather than this newly proposed course (ECON 8877), especially since methodological topics are not covered extensively in the other two courses. Of course, cancellation decisions would be clearly be made on a year-to-year basis depending on the current availability of instructors.

\bigskip


Paul J. Healy

\end{document}
